\documentclass[11pt]{article}
\usepackage[margin=1in]{geometry}
\usepackage{amsmath}
\usepackage[hidelinks]{hyperref}
\title{Swimming: Distance and Sets}
\author{Liam McGrath}
\date{4 October 2026}
\begin{document}
\maketitle

\section{Overview}
Swimming has long been used as a low-impact warm-up before weight training, and as a standalone workout on rest days. Every stroke is done as equally as possible, to balance the use of each muscle group.

\section{Set training, next}
When a proper water-safe watch is bought, set data will be recorded for $5 \times 50$ m, $10 \times 50$ m and $5 \times 100$ m. Each repeat is to be completed within a preset time window (for example, 48 s for each 50 m), and the remaining time is used for rest.

\section{The high-school rivalry}
This kind of training was used often in high school, where Liam and Owen were rivals who competed in the 50 m freestyle and the 50 m butterfly. Their best times are below.

\begin{table}[htbp]
\centering
\begin{tabular}{lll}
\hline
Event & Liam & Owen \\
\hline
50 m freestyle & 27.09 s & 26.99 s \\
50 m butterfly & 32.60 s & 33.80 s \\
\hline
\end{tabular}
\caption{High-school best times, as recalled.}
\end{table}

\section{The current log}
Total distances have been logged since 21 September 2026. Each session so far is a set of 400 m swims in the order fly, back, breast, free. The Swimming page of the website charts every session and the running total, with the log as JSON and CSV downloads.

\end{document}
